\chapter{Introduction}
\label{sec:introduction}

Agentic Artificial Intelligence (Agentic AI) denotes a class of AI systems engineered
to function as autonomous, goal-oriented agents. These systems are capable of perceiving
their environment, reasoning about defined objectives, and executing actions with minimal
human intervention. Agentic AI represents an evolution of traditional rule-based control
systems, improving upon the latter by integrating advanced decision-making components,
such as reinforcement learning algorithms or large language models (LLMs), that
serve as the central reasoning engine. These core models are complemented by additional
architectural components, including persistent memory modules, planning algorithms,
specialized tool interfaces, and feedback mechanisms. Together, these elements enable
structured, multi-step task execution and adaptive behavior in dynamic
environments~\cite{agentic2025acharya}.
Owing to these capabilities, Agentic AI systems are applicable across a wide range of
domains, including healthcare, cybersecurity, autonomous software development, industrial
control systems, and process automation.

In the field of information technology, such solutions offer a compelling approach to
managing systems that are becoming increasingly large, dynamic, and complex in order to
meet evolving market demands. In this context, the \gls{AIOps} paradigm~\cite{GartnerAIOpsDefinition}
represents a shift from traditional operations, primarily driven by human expertise
and manual effort, toward partially or fully autonomous processes enabled by AI.

This transition presents both opportunities and challenges. On the one hand, the
integration of AI models into operational workflows can streamline processes,
significantly accelerate decision-making, reduce operational costs, and mitigate the
impact of human error. On the other hand, the practical adoption of AIOps systems
entails addressing a multitude of critical challenges, including the identification
of appropriate use cases, the observability and quality of contextual data required
to adapt and operate AI models, interfacing said models with their target environments,
the deployment of large-scale models in resource-constrained environments, the
scalability of solutions to massive infrastructures, and the limited transparency
and interpretability of AI-driven decision-making mechanisms.
Additionally, granting a high degree of autonomy to AI systems introduces significant
safety concerns, particularly in critical scenarios where incorrect decisions may
result in service degradation or disruption for end users. Consequently, the integration
of AI agents must be accompanied by robust safeguard mechanisms and, where appropriate,
human oversight, which remains a desirable requirement in current deployments.

Specifically in the context of cloud systems, \gls{LLM}-based Agentic AI possesses a
great deal of unexpressed potential, both as a helper for human operators and as
fully autonomous entities capable of running, orchestrating, and troubleshooting
operations. Agents may prove useful for detecting anomalies, predicting faults,
performing root cause analysis, and allocating resources~\cite{cheng2023ai}.
Recently, efforts have been expended in defining guidelines and design principles,
paving the road towards the development of novel solutions in this
field~\cite{shetty2024building}; however, many research avenues remain unexplored.

\section{Problem Statement}
\label{sec:problem}

Kubernetes has become the de facto standard for orchestrating containerized
workloads in production cloud environments. Its built-in autoscaling mechanisms—the
\gls{HPA} and the \gls{VPA}—address resource management by adjusting replica counts
and resource limits in response to observed metrics. Despite their widespread adoption,
both mechanisms share a fundamental limitation: they are \textit{reactive}. The HPA
waits for \gls{CPU} utilization to cross a threshold before scaling; the VPA waits
for a stable observation window before updating resource requests. In both cases,
provisioning lag allows SLA violations to accumulate before additional capacity
arrives.

Three further problems motivate this work:

\begin{enumerate}
  \item \textbf{No forecasting.} Neither HPA nor VPA considers predicted future load.
  An approaching traffic burst is invisible until it materializes as metric pressure.

  \item \textbf{Single-criterion optimization.} HPA optimizes for a single metric
  (typically CPU). Real operational trade-offs—cost vs.\ performance vs.\ stability—are
  not considered.

  \item \textbf{No explainability.} Neither mechanism produces a human-readable
  explanation for its actions. An operator observing an unexpected scale event cannot
  determine why it occurred without inspecting raw metric history and control-loop
  logs. This opacity is a barrier in regulated environments and in organizations
  seeking auditability of infrastructure cost.
\end{enumerate}

\section{Proposed Approach and Contributions}
\label{sec:contributions}

In this setting, we introduce \textbf{Ai4k8s}, an agentic AI platform engineered to
support safe, intelligent Kubernetes operations. Ai4k8s leverages the capabilities
of \glspl{LLM} to make dynamic, context-aware autoscaling decisions and provides a
full decision trace for every action it takes. Ai4k8s implements a three-layer pipeline:

\begin{enumerate}
  \item \textbf{Enforcement layer} — applies hard safety rules to block physically
  impossible or operationally unsafe scaling actions before any model is consulted.

  \item \textbf{LLM reasoning layer} — a quantized local language model
  (\textit{qwen3:0.6b}, 522~MB) reads live cluster metrics, a time-series forecast,
  and deployment metadata to produce a plain-English scaling recommendation. A
  graceful cascade (local LLM $\to$ cloud LLM $\to$ regex fallback) ensures
  availability under resource contention.

  \item \textbf{\gls{MCDA} optimization layer} — a TOPSIS optimizer scores all
  candidate actions across five criteria (cost, performance, stability, forecast
  alignment, response time) and confirms or overrides the LLM recommendation.
\end{enumerate}

The principal contributions of this thesis are:

\begin{itemize}
  \item A novel three-layer autoscaling architecture that combines deterministic
  safety rules, \gls{LLM} reasoning, and multi-criteria optimization in a single
  production-grade pipeline.

  \item Demonstration that a sub-gigabyte on-premises LLM is sufficient to drive
  Kubernetes autoscaling at 1.7~s latency on commodity hardware, with no cloud
  dependency at idle.

  \item A 52\% reduction in compute cost proxy compared to native HPA, with full
  decision explainability, validated on the muBench microservice benchmark.

  \item An open-source implementation available at \url{https://ai4k8s.online} with
  a live web dashboard, MCP tool server, and structured evaluation harness.
\end{itemize}

\section{Thesis Structure}
\label{sec:structure}

The remainder of this thesis is organized as follows:

\begin{itemize}
  \item \textbf{Chapter~\ref{ch:background}} introduces the background on Kubernetes
  autoscaling, the AIOps paradigm, LLMs as reasoning engines, and multi-criteria
  decision analysis. It surveys related work and positions Ai4k8s in the literature.

  \item \textbf{Chapter~\ref{ch:design}} describes the Ai4k8s system architecture
  in detail, covering the Ai4k8s decision pipeline, the local LLM deployment
  strategy, the MCDA optimizer, and the predictive monitoring subsystem.

  \item \textbf{Chapter~\ref{ch:implementation}} discusses the implementation,
  focusing on the engineering challenges encountered—LLM cascade design, Ollama
  integration, TOPSIS weight calibration, and the parallel background autoscaling
  loop—and the solutions developed.

  \item \textbf{Chapter~\ref{ch:evaluation}} presents the experimental setup and
  results: bottleneck analysis with seq\_len tuning, predictive scale-up
  demonstration, HPA/VPA mode selection accuracy, and the structured three-way
  comparison (HPA vs.\ VPA vs.\ Ai4k8s) across multiple load levels.

  \item \textbf{Chapter~\ref{ch:conclusions}} summarizes the findings, discusses
  limitations, and outlines directions for future work.
\end{itemize}